% ============================================================
%  Paste this entire file into Overleaf → New Project → Blank
%  Compiler: pdfLaTeX  |  Main file: research_paper.tex
% ============================================================
\documentclass[11pt,a4paper]{article}

% ── Packages ─────────────────────────────────────────────────
\usepackage[margin=2.5cm]{geometry}
\usepackage[T1]{fontenc}
\usepackage[utf8]{inputenc}
\usepackage{lmodern}
\usepackage{microtype}
\usepackage{times}
\usepackage{amsmath,amssymb}
\usepackage{graphicx}
\usepackage{booktabs}
\usepackage{multirow}
\usepackage{array}
\usepackage{xcolor}
\usepackage{url}
\usepackage{hyperref}
\usepackage{enumitem}
\usepackage{caption}
\usepackage{subcaption}
\usepackage{float}
\usepackage{parskip}
\usepackage{setspace}
\usepackage{fancyhdr}
\usepackage{abstract}
\usepackage{titlesec}
\usepackage[numbers,sort&compress]{natbib}

% ── Colours ───────────────────────────────────────────────────
\definecolor{darkblue}{HTML}{1a237e}
\definecolor{accent}{HTML}{0d47a1}
\definecolor{lightgray}{HTML}{f5f5f5}
\definecolor{tableheader}{HTML}{e8eaf6}

% ── Hyperref setup ────────────────────────────────────────────
\hypersetup{
  colorlinks=true,
  linkcolor=accent,
  citecolor=accent,
  urlcolor=accent,
  pdftitle={Quadruplet Sentence Alignment for Legal Machine Translation},
  pdfauthor={Darsh Iyer}
}

% ── Section formatting ────────────────────────────────────────
\titleformat{\section}{\large\bfseries\color{darkblue}}{{\thesection.}}{0.5em}{}
\titleformat{\subsection}{\normalsize\bfseries\color{accent}}{{\thesubsection}}{0.5em}{}
\titleformat{\subsubsection}{\normalsize\itshape}{\thesubsubsection}{0.5em}{}

% ── Header / footer ───────────────────────────────────────────
\pagestyle{fancy}
\fancyhf{}
\renewcommand{\headrulewidth}{0.4pt}
\fancyhead[L]{\small\textit{Quadruplet Alignment for Legal MT — English, Marathi, Kannada, Hindi}}
\fancyhead[R]{\small\thepage}

% ── Table helpers ─────────────────────────────────────────────
\newcolumntype{C}[1]{>{\centering\arraybackslash}p{#1}}
\newcolumntype{L}[1]{>{\raggedright\arraybackslash}p{#1}}

\setlength{\tabcolsep}{6pt}
\renewcommand{\arraystretch}{1.25}

% ── Macros ────────────────────────────────────────────────────
\newcommand{\mr}{\textsc{mr}}
\newcommand{\kn}{\textsc{kn}}
\newcommand{\en}{\textsc{en}}
\newcommand{\hi}{\textsc{hi}}
\newcommand{\bleu}{\textsc{bleu}}
\newcommand{\chrf}{\textsc{chr}F}
\newcommand{\bert}{\textsc{bert}Score}
\newcommand{\acc}{Acc@1}
\newcommand{\mrr}{\textsc{mrr}}
\newcommand{\nllb}{\textsc{nllb}-200}

% ─────────────────────────────────────────────────────────────
\begin{document}

% ── Title ─────────────────────────────────────────────────────
\begin{titlepage}
  \centering
  \vspace*{1.5cm}

  {\color{darkblue}\rule{\linewidth}{2pt}}
  \vspace{0.4cm}

  {\LARGE\bfseries\color{darkblue}
    Quadruplet Sentence Alignment for\\[0.3em]
    Low-Resource Legal Machine Translation:\\[0.3em]
    English--Marathi--Kannada--Hindi
  }

  \vspace{0.4cm}
  {\color{darkblue}\rule{\linewidth}{0.8pt}}

  \vspace{1.2cm}

  {\large Darsh Iyer}\\[0.2em]
  {\normalsize Department of Computer Science}\\[0.1em]
  {\normalsize \textit{Research Intern, NLP \& Machine Translation Group}}\\[0.5em]
  {\small \href{mailto:darsh@example.com}{darsh@example.com}}

  \vspace{1cm}
  {\small Supervised by Vikram Sir}

  \vspace{1.5cm}
  {\small April 2026}

  \vfill
  \colorbox{tableheader}{%
    \begin{minipage}{0.85\linewidth}
      \vspace{0.3cm}
      \centering
      \small\textbf{Keywords:} Multilingual Machine Translation,
      Sentence Alignment, Low-Resource NLP, Indic Languages,
      Quadruplet Training, Legal Domain MT, English--Marathi--Kannada
      \vspace{0.3cm}
    \end{minipage}
  }
\end{titlepage}

% ── Abstract ──────────────────────────────────────────────────
\newpage
\begin{abstract}
\noindent
Direct machine translation between Marathi and Kannada is severely
under-resourced: no standardised parallel benchmark exists, and
the two languages share neither script nor close linguistic ancestry.
We address this gap through a \emph{pivot-based quadruplet alignment}
framework that leverages English as an intermediary language and
Hindi as an additional Indic anchor to improve cross-lingual
embedding coherence.

We construct a corpus of \textbf{128,460 sentence-level quadruplets}
(English, Marathi, Kannada, Hindi) drawn from four parallel sources:
Samanantar, MILPaC (legal domain), PMIndia (government speeches),
and OPUS-100.
We fine-tune a multilingual sentence transformer
(\texttt{paraphrase-multilingual-MiniLM-L12-v2}) under two
conditions --- triplet (\en+\mr+\kn) and quadruplet
(\en+\mr+\kn+\hi) --- and evaluate retrieval accuracy
across all six language-pair directions on a test pool of 25,692
sentences.

Our quadruplet model achieves \acc{} of \textbf{46.3\%} for
\en{}$\leftrightarrow$\mr{} and \textbf{40.7\%} for
\en{}$\leftrightarrow$\kn{}, compared to a 7.8\% and 41.0\%
baseline respectively, demonstrating substantial gains especially
for Kannada.
For the direct Marathi$\leftrightarrow$Kannada direction, the
triplet model achieves 30.2\% \acc{}, a 4.5$\times$ improvement
over the 6.8\% baseline.
We further evaluate pivot translation quality using \nllb{},
reporting \bleu{}/\chrf{}/\bert{}Score, and analyse English
structural bias through centroid cosine similarity.
All code, models, and datasets are publicly released.
\end{abstract}

\vspace{0.5cm}
\tableofcontents
\newpage

% ── 1. INTRODUCTION ───────────────────────────────────────────
\section{Introduction}
\label{sec:intro}

India is home to 22 constitutionally recognised languages and over
1,600 dialects, yet the NLP infrastructure for most of these
languages lags decades behind English.
Marathi (spoken by $\sim$83 million people) and Kannada
($\sim$44 million) are two of the six major Dravidian and
Indo-Aryan languages of the southern and western regions, and their
legal corpora --- court judgments, Acts of Parliament, government
notifications --- are produced simultaneously in multiple languages
by constitutional mandate.
Despite this, no production-quality machine translation system
exists for the direct Marathi$\leftrightarrow$Kannada language pair.

\subsection{Problem Statement}
The central challenge is \textbf{data scarcity}: there is no publicly
available direct Marathi--Kannada parallel corpus of meaningful size.
Existing multilingual systems such as IndicTrans2~\cite{gala2023indictrans2}
and NLLB-200~\cite{nllb2022} support both languages but only via
English as an implicit pivot, without exposing sentence-level
alignment quality or domain-specific performance.

Furthermore, the legal domain imposes stricter requirements than
general news translation: terminology must be precise, clause
boundaries must be preserved, and cross-reference integrity must
be maintained.
Errors in legal MT can lead to misinterpretation of statutes,
a concern increasingly relevant as Indian courts digitise their records.

\subsection{Contributions}
This paper makes the following contributions:

\begin{enumerate}[leftmargin=*, label=\textbf{C\arabic*.}]
  \item \textbf{Quadruplet corpus:} We compile and release a
    corpus of 128,460 sentence-level quadruplets
    (EN + MR + KN + HI) from four data sources, with
    full preprocessing and alignment pipeline.

  \item \textbf{Comparative embedding study:}
    We systematically compare baseline (no fine-tuning),
    triplet (\en+\mr+\kn), and quadruplet (\en+\mr+\kn+\hi)
    training regimes across six language-pair directions and
    a 25,692-sentence retrieval pool.

  \item \textbf{Bias analysis:}
    We introduce a centroid cosine similarity metric to quantify
    English structural bias in the embedding space and show how
    adding Hindi affects Indic language clustering.

  \item \textbf{Legal domain translation evaluation:}
    We report \bleu{}, \chrf{}, and \bert{}Score for
    Marathi$\rightarrow$Kannada, Kannada$\rightarrow$Marathi, and
    English$\rightarrow$\{Marathi, Kannada\} directions using the
    \nllb{} pivot pipeline.

  \item \textbf{Open prototype:}
    We release a FastAPI backend and interactive dashboard
    enabling PDF upload, language detection, and pivot translation,
    available at \url{https://github.com/darshiyer/internship}.
\end{enumerate}

% ── 2. RELATED WORK ───────────────────────────────────────────
\section{Related Work}
\label{sec:related}

\subsection{Multilingual Sentence Embeddings}
The field of language-agnostic sentence representations was advanced
significantly by LASER~\cite{artetxe2019massively}, which trained
a single LSTM encoder across 93 languages using a shared BPE
vocabulary. LaBSE~\cite{feng2022labse} subsequently demonstrated
that a BERT-based architecture with dual-encoder contrastive training
could outperform LASER on bitext mining tasks, including for
morphologically rich languages.
Reimers and Gurevych~\cite{reimers2019sentencebert} introduced
Sentence-BERT (SBERT), extending BERT with siamese networks and
cosine similarity losses; their multilingual MiniLM variant
(\texttt{paraphrase-multilingual-MiniLM-L12-v2}) is our base model.

\subsection{Pivot-Based Machine Translation}
Pivot translation -- routing through a high-resource intermediate
language -- is a well-established technique for low-resource
pairs~\cite{utiyama2007pivot,wu2007pivot}.
\citet{paul2009pivotmt} showed that for Asian language pairs with
limited direct data, English pivot can match or exceed direct
models when the pivot language is richly resourced.
The key limitation is \emph{error propagation}: errors in the
source$\rightarrow$pivot step compound with errors in the
pivot$\rightarrow$target step~\cite{babych2007pivotmt}.
Our round-trip consistency experiment (Section~\ref{sec:results})
directly quantifies this effect.

\subsection{Indian Language NLP}
Samanantar~\cite{ramesh2021samanantar} is the largest publicly
available parallel corpus for Indian language pairs, providing
English-aligned data for 11 Indic languages including Marathi and
Kannada. We use 500,000 sentences per language pair from this source.
MILPaC~\cite{milpac2022} provides a legal-domain parallel corpus for
Indian languages, which we incorporate for domain adaptation.
IndicTrans2~\cite{gala2023indictrans2} achieves state-of-the-art
translation quality for 22 Indian languages; we use it as our
Indic$\rightarrow$English translation backbone.
\citet{nllb2022} released NLLB-200, a 200-language model that we
use for English$\rightarrow$Indic translation.

\subsection{Legal Machine Translation}
Legal MT presents distinct challenges~\cite{prieto2015legal}:
dense nominalisation, passive constructions, cross-reference
clauses, and archaic terminology.
\citet{koehn2017challenges} identified domain mismatch as the
primary cause of quality degradation in specialised MT.
Our inclusion of MILPaC and PMIndia (government/policy domain)
addresses this partially; Section~\ref{sec:limitations} discusses
remaining domain gaps.

% ── 3. DATASET ────────────────────────────────────────────────
\section{Dataset Construction}
\label{sec:data}

\subsection{Data Sources}
We aggregate parallel sentence pairs from four sources,
summarised in Table~\ref{tab:datasources}.

\begin{table}[H]
\centering
\caption{Parallel corpus sources used in this work.
  All pairs are sentence-level with English as one side.}
\label{tab:datasources}
\small
\renewcommand{\arraystretch}{1.3}
\begin{tabular}{L{2.8cm} C{1.4cm} C{1.4cm} C{1.4cm} L{3.6cm}}
\toprule
\textbf{Source} & \textbf{EN--MR} & \textbf{EN--KN} & \textbf{EN--HI} & \textbf{Domain} \\
\midrule
Samanantar~\cite{ramesh2021samanantar}
  & 500k & 500k & 500k & General (news, web) \\
MILPaC~\cite{milpac2022}
  & incl. & incl. & incl. & Legal (acts, judgments) \\
PMIndia (this work)
  & 36,110 & 35,178 & 56,792 & Govt.\ speeches \\
OPUS-100~\cite{zhang2020opus100}
  & 22,432 & 10,608 & 75,611 & Mixed \\
\midrule
\textbf{Total (after dedup)}
  & \textbf{553,939} & \textbf{542,876} & \textbf{499,979} & --- \\
\bottomrule
\end{tabular}
\end{table}

\subsection{Quadruplet Alignment}
\label{sec:alignment}

A quadruplet $Q_i = (\text{en}_i, \text{mr}_i, \text{kn}_i, \text{hi}_i)$
requires a sentence that appears in \emph{all four} language-specific
corpora.
Since the EN--MR and EN--KN source documents originate from
different publications, direct index matching is insufficient.
We apply a \textbf{three-tier English key normalisation} pipeline:

\begin{enumerate}[leftmargin=*, label=\arabic*.]
  \item \textit{Exact match}: UTF-8 normalised (NFC) lowercase string equality.
  \item \textit{Whitespace-normalised match}: Collapse all whitespace, strip
    leading/trailing punctuation.
  \item \textit{Stemmed match}: Apply lightweight suffix stripping
    (remove trailing \texttt{-ing}, \texttt{-ed}, \texttt{-s}).
\end{enumerate}

Matches at tier 1 receive confidence score 1.0; tiers 2 and 3
receive 0.9 and 0.8 respectively.
Only pairs with matching English keys across at least the MR and KN
corpora are retained as quadruplets; the HI side is optional
(creating \emph{full quadruplets} vs.\ \emph{triplets}).

\subsection{Dataset Statistics}
Table~\ref{tab:quadstats} summarises the final aligned dataset.

\begin{table}[H]
\centering
\caption{Final quadruplet dataset statistics after alignment and deduplication.}
\label{tab:quadstats}
\begin{tabular}{lrr}
\toprule
\textbf{Split / Type} & \textbf{Count} & \textbf{Percentage} \\
\midrule
Full quadruplets (EN+MR+KN+HI) & 28,794  & 22.4\% \\
Triplets (EN+MR+KN, no HI)     & 99,666  & 77.6\% \\
\midrule
\textbf{Total quadruplets}      & \textbf{128,460} & 100\% \\
\midrule
Train split & 89,922 & 70\% \\
Development split & 12,846 & 10\% \\
Test split  & 25,692 & 20\% \\
\bottomrule
\end{tabular}
\end{table}

\subsection{Data Cleaning}
All source corpora were cleaned as follows:
\begin{itemize}[leftmargin=*]
  \item Unicode NFC normalisation; removal of zero-width characters.
  \item Deduplication on normalised English key.
  \item Length filtering: $10 \leq |\text{tokens}| \leq 150$ per sentence.
  \item Language verification using Unicode block analysis:
    Devanagari (U+0900--U+097F) for MR/HI;
    Kannada (U+0C80--U+0CFF) for KN.
  \item Length-ratio filtering: $0.3 \leq |s_\text{en}|/|s_\text{tgt}| \leq 4.0$
    (character counts; Indic scripts are denser than Latin).
\end{itemize}

% ── 4. METHODOLOGY ────────────────────────────────────────────
\section{Methodology}
\label{sec:method}

\subsection{Sentence Embedding Model}

We use \texttt{paraphrase-multilingual-MiniLM-L12-v2}
as our base encoder.
This 12-layer distilled BERT model supports 50+ languages,
including Marathi, Kannada, and Hindi, and produces
384-dimensional sentence embeddings.
We compare three training conditions:

\begin{description}[leftmargin=*, font=\bfseries]
  \item[Baseline] The pre-trained model with no fine-tuning.
    Used as the zero-shot reference.
  \item[Triplet] Fine-tuned on \en+\mr+\kn sentence triplets
    using cosine similarity loss. The model sees all three
    language variants of the same sentence.
  \item[Quadruplet] Fine-tuned on \en+\mr+\kn+\hi quadruplets.
    Hindi provides an additional Indic anchor, hypothesised to
    reduce English structural bias (Section~\ref{sec:bias}).
\end{description}

\subsection{Training Details}
\label{sec:training}

For both fine-tuned models we use the following hyperparameters,
selected to balance quality against CPU training time on Apple
Silicon M-series hardware:

\begin{table}[H]
\centering
\caption{Training hyperparameters.}
\label{tab:hparams}
\begin{tabular}{ll}
\toprule
\textbf{Hyperparameter} & \textbf{Value} \\
\midrule
Base model & \texttt{paraphrase-multilingual-MiniLM-L12-v2} \\
Loss function & CosineSimilarityLoss \\
Training pairs & 8,000 (sampled from 179,844 / 330,369 available) \\
Batch size & 32 \\
Epochs & 3 \\
Warmup steps & 100 \\
Device & CPU (MPS disabled; Apple Silicon OOM with NLLB-600M) \\
Training time & $\sim$35 min per model \\
\bottomrule
\end{tabular}
\end{table}

All pairs are constructed by taking every valid language combination
within each quadruplet (e.g., EN--MR, EN--KN, MR--KN for triplets;
additionally EN--HI, MR--HI, KN--HI for quadruplets), yielding
$\binom{3}{2}=3$ and $\binom{4}{2}=6$ pair types respectively.

\subsection{Translation Pipeline}
\label{sec:translation}

We implement an English-pivot translation chain:

\[
\text{Marathi} \xrightarrow{\textsc{IndicTrans2}} \text{English}
\xrightarrow{\textsc{NLLB-200}} \text{Kannada}
\]

\begin{itemize}[leftmargin=*]
  \item \textbf{Indic$\rightarrow$English:} \texttt{ai4bharat/indictrans2-indic-en-1B}.
    Supports all 22 Indian languages scheduled in the Constitution.
  \item \textbf{English$\rightarrow$Indic (general):}
    \texttt{facebook/nllb-200-distilled-600M}.
    Selected for balance between quality and inference speed.
  \item \textbf{English$\rightarrow$Indic (legal):}
    \texttt{law-ai/InLegalTrans-En2Indic-1B}, fine-tuned on MILPaC.
\end{itemize}

Translation evaluation is performed on a held-out set of 100
sentence pairs (MR+KN+EN triples from the test split),
using a GPU (T4) on Google Colab to avoid Apple Silicon
out-of-memory errors with the NLLB-600M model.

\subsection{Evaluation Metrics}

\paragraph{Retrieval (Embedding Quality)}
For each language pair $(L_1, L_2)$, we embed all test sentences
in both languages and compute pairwise cosine similarity.
\acc{} (Accuracy at 1) measures the fraction of source sentences
whose top-1 nearest neighbour in the target language is the correct
translation.
Mean Reciprocal Rank (\mrr{}) provides a rank-averaged measure.
Both metrics are computed over the full 25,692-sentence test pool.

\paragraph{Translation Quality}
\begin{itemize}[leftmargin=*]
  \item \bleu{}~\cite{papineni2002bleu}: Standard $n$-gram precision
    metric with brevity penalty.
  \item \chrf{}~\cite{popovic2015chrf}: Character $n$-gram F-score,
    more robust than \bleu{} for morphologically rich languages.
  \item \bert{}Score~\cite{zhang2020bertscore}: Precision/Recall/F1
    computed from contextual embeddings of hypothesis and reference.
  \item \textbf{Round-trip consistency:} Cosine similarity between
    the original English sentence and the embedding of its
    Marathi back-translation.
\end{itemize}

% ── 5. RESULTS ────────────────────────────────────────────────
\section{Experiments and Results}
\label{sec:results}

\subsection{Embedding Retrieval Results}

Table~\ref{tab:acc1} reports \acc{} and \mrr{} for all three
models across five language-pair directions on the 25,692-sentence
test pool.

\begin{table}[H]
\centering
\caption{Retrieval performance (\acc{} and \mrr{}) on the 25,692-sentence test pool.
  Bold indicates the best fine-tuned result per row.
  $n$ = pool size per language pair.}
\label{tab:acc1}
\small
\begin{tabular}{l C{1.5cm} C{1.5cm} C{1.5cm} C{1.5cm} C{1.5cm} C{1.5cm}}
\toprule
& \multicolumn{2}{c}{\textbf{Baseline}} & \multicolumn{2}{c}{\textbf{Triplet}} & \multicolumn{2}{c}{\textbf{Quadruplet}} \\
\cmidrule(lr){2-3}\cmidrule(lr){4-5}\cmidrule(lr){6-7}
\textbf{Pair} & \acc & \mrr & \acc & \mrr & \acc & \mrr \\
\midrule
EN$\leftrightarrow$MR & 41.0 & .500 & 42.7 & .525 & \textbf{46.3} & \textbf{.560} \\
EN$\leftrightarrow$KN &  7.8 & .115 & 39.2 & .495 & \textbf{40.7} & \textbf{.512} \\
EN$\leftrightarrow$HI & 50.5 & .639 & 48.4 & .614 & \textbf{50.3} & \textbf{.635} \\
MR$\leftrightarrow$KN &  6.8 & .099 & \textbf{30.2} & \textbf{.393} & 29.6 & .388 \\
MR$\leftrightarrow$HI & 32.5 & .441 & \textbf{35.7} & \textbf{.483} & 34.6 & .466 \\
\bottomrule
\multicolumn{7}{l}{\small\textit{Note: EN$\leftrightarrow$HI and MR$\leftrightarrow$HI evaluated on 5,680 full-quadruplet test sentences.}}
\end{tabular}
\end{table}

\paragraph{Key findings.}
\begin{itemize}[leftmargin=*]
  \item \textbf{EN$\leftrightarrow$KN improvement is dramatic.}
    The baseline model achieves only 7.8\% \acc{} for
    English$\leftrightarrow$Kannada, suggesting that the
    pre-trained model has very limited Kannada representation.
    Both fine-tuned models raise this to $\sim$40\%, a
    \textbf{5.2$\times$ improvement}.
  \item \textbf{Quadruplet beats triplet for EN-centric pairs.}
    For EN$\leftrightarrow$MR (+3.6\%) and EN$\leftrightarrow$KN
    (+1.5\%), the quadruplet model consistently outperforms the
    triplet, suggesting that Hindi provides useful cross-lingual
    signal beyond Marathi and Kannada alone.
  \item \textbf{Triplet wins for direct MR$\leftrightarrow$KN.}
    The triplet model achieves 30.2\% vs.\ 29.6\% for quadruplet
    on the direct Marathi--Kannada pair. This is within margin of
    error and likely reflects the trade-off between
    MR$\leftrightarrow$KN pairs available in each training regime.
  \item \textbf{Retrieval pool size matters.}
    The 25,692-sentence pool is $\sim$45$\times$ larger than the
    563-sentence pools used in initial experiments, making the
    task substantially harder and the reported numbers more
    trustworthy as real-world estimates.
\end{itemize}

\subsection{Translation Quality Results}

Table~\ref{tab:translation} reports translation quality on the
100-sentence held-out test set evaluated on Google Colab (T4 GPU).

\begin{table}[H]
\centering
\caption{Pivot translation quality using \nllb{} (distilled-600M).
  All experiments use the legal domain model (InLegalTrans) for EN$\rightarrow$Indic.}
\label{tab:translation}
\begin{tabular}{l C{1.8cm} C{1.8cm} C{1.8cm} C{3.0cm}}
\toprule
\textbf{Direction} & \textbf{\bleu} & \textbf{\chrf} & \textbf{\bert{}F1} & \textbf{Note} \\
\midrule
MR $\rightarrow$ EN $\rightarrow$ KN & 2.77  & 32.50 & 0.7865 & Pivot: EN bridge \\
KN $\rightarrow$ EN $\rightarrow$ MR & 7.44  & 32.41 & 0.7855 & Pivot: EN bridge \\
EN $\rightarrow$ KN (direct)         & 3.04  & 36.22 & 0.8076 & Direct EN input \\
EN $\rightarrow$ MR $\rightarrow$ EN & --- & --- & cos = 0.898 & Round-trip \\
\bottomrule
\end{tabular}
\end{table}

\paragraph{Interpreting the scores.}
\bleu{} scores below 10 are common for distant language pairs
and short test sets; they do not indicate a broken system.
The more informative metrics are \chrf{} ($\sim$32--36, indicating
substantial character-level overlap) and \bert{}Score F1
($\sim$0.786--0.808), which measures semantic rather than lexical
similarity and is more appropriate for Indic languages with
rich morphology.
The round-trip consistency score of 0.898 cosine similarity
confirms that semantic content is largely preserved through the
English pivot.

\paragraph{KN$\rightarrow$MR outperforms MR$\rightarrow$KN.}
The BLEU gap (7.44 vs.\ 2.77) likely reflects the relative
abundance of EN--MR training data for the target-side model.
Kannada is a morphologically richer language (agglutinative)
which inflates reference sentence length relative to output,
penalising the brevity-sensitive \bleu{} metric.

% ── 6. BIAS ANALYSIS ──────────────────────────────────────────
\section{Bias Analysis: English Structural Bias}
\label{sec:bias}

A well-known failure mode of multilingual models trained
predominantly on English data is \emph{English structural bias}:
sentences in different languages are embedded close to their
English translations rather than to each other, causing
the embedding space to be organised around English grammatical
structure rather than semantic content~\cite{pires2019multibert}.

\subsection{Centroid Cosine Similarity Metric}

We quantify this effect using the \textbf{centroid cosine
similarity}: for a language pair $(L_1, L_2)$, we compute
the cosine similarity between the centroid of all $L_1$
embeddings and the centroid of all $L_2$ embeddings.
A value near 1.0 indicates the two language embedding clouds
are collapsed together; a value near 0 indicates full separation.
For an unbiased multilingual model, we would expect
centroid similarities to be uniformly high across all language pairs.

\begin{table}[H]
\centering
\caption{Centroid cosine similarity between language pair embedding clouds
  (computed on 25,692 test sentences).
  Higher = less geometric separation between languages.}
\label{tab:centroid}
\begin{tabular}{l C{2.2cm} C{2.2cm} C{2.2cm}}
\toprule
\textbf{Language Pair} & \textbf{Baseline} & \textbf{Triplet} & \textbf{Quadruplet} \\
\midrule
EN $\leftrightarrow$ MR & 0.9108 & 0.9995 & 0.9990 \\
EN $\leftrightarrow$ KN & 0.5337 & 0.9992 & 0.9988 \\
EN $\leftrightarrow$ HI & 0.8703 & 0.9987 & 0.9982 \\
MR $\leftrightarrow$ KN & 0.6338 & 0.9998 & 0.9998 \\
\bottomrule
\end{tabular}
\end{table}

\subsection{Findings}

\paragraph{The baseline model shows strong bias against Kannada.}
The EN$\leftrightarrow$KN centroid similarity of 0.534 in the
baseline model is strikingly lower than EN$\leftrightarrow$MR
(0.911) and EN$\leftrightarrow$HI (0.870), indicating that
the base model's Kannada representations are geometrically
distant from both English and Devanagari-script languages.
This explains the extremely low baseline \acc{} of 7.8\%
for EN$\leftrightarrow$KN.

\paragraph{Fine-tuning collapses the embedding space.}
Both triplet and quadruplet models produce centroid similarities
of $>$0.998 across all language pairs, indicating that
fine-tuning has pulled all four language embedding clouds
into a tight cluster.
This is expected given the cosine similarity training loss,
which directly optimises for this alignment.

\paragraph{Does Hindi reduce English bias?}
The centroid values are similar between triplet and quadruplet
($\leq$0.001 difference), suggesting that within this metric,
Hindi's addition provides marginal geometric benefit.
However, the retrieval improvements in Table~\ref{tab:acc1}
indicate that Hindi does contribute useful signal for the EN-centric
pairs, likely by providing additional contextual grounding
rather than changing the geometric distribution.

\paragraph{Limitation of this metric.}
High centroid similarity in fine-tuned models reflects
\emph{language collapse} rather than \emph{semantic alignment}:
all embeddings may be concentrated in a small region regardless
of content.
A more informative measure would be \emph{intra-language spread}
(average pairwise distance within a language's embedding cloud),
which we leave as future work.

% ── 7. LIMITATIONS ────────────────────────────────────────────
\section{Limitations}
\label{sec:limitations}

We report limitations explicitly, following the honest documentation
strategy recommended by our Systematic Literature Review framework.

\subsection{Dataset Limitations}
\begin{itemize}[leftmargin=*]
  \item \textbf{No standardised MR$\leftrightarrow$KN benchmark.}
    There is no publicly available, human-annotated evaluation
    benchmark for direct Marathi--Kannada translation.
    Our test set is constructed automatically from the alignment
    pipeline and may contain noise.

  \item \textbf{Domain imbalance.}
    The majority of our quadruplets originate from Samanantar
    (general news/web domain) and PMIndia (government speeches).
    Purely legal text (from MILPaC) forms a small proportion.
    Legal-domain performance may be overestimated by mixing with
    general-domain evaluation.

  \item \textbf{Pivot-based alignment.}
    All MR--KN sentence pairs are aligned via English keys,
    not by direct comparison of Marathi and Kannada text.
    Sentences with equivalent meaning but different English
    translations may be incorrectly separated.
\end{itemize}

\subsection{Modelling Limitations}
\begin{itemize}[leftmargin=*]
  \item \textbf{Small training set.}
    We train on only 8,000 sentence pairs per model due to
    CPU time constraints, from a pool of $>$179,000 available.
    Models trained on the full dataset would likely show
    further improvement.

  \item \textbf{Embedding collapse.}
    The centroid similarity results indicate that
    CosineSimilarityLoss may over-constrain the embedding
    space, pulling all language representations into
    a tight cluster.
    Contrastive losses (e.g., InfoNCE, MultipleNegativesRankingLoss)
    may mitigate this.

  \item \textbf{No end-to-end fine-tuning of the translation model.}
    We use the pre-trained NLLB-200 and IndicTrans2 models
    without domain adaptation on legal text.
    Fine-tuning these on our legal corpus would likely
    improve \bleu{}/\chrf{} scores substantially.
\end{itemize}

\subsection{Evaluation Limitations}
\begin{itemize}[leftmargin=*]
  \item \textbf{Absence of human evaluation.}
    All metrics (\bleu{}, \chrf{}, \bert{}Score) are automatic.
    Human evaluation for morphology, word order, named entity
    preservation, and omission errors would provide a more
    complete picture.

  \item \textbf{\bleu{} not suitable for morphologically rich
    languages.}
    Both Marathi and Kannada are morphologically rich, with
    agglutinative morphology in Kannada and fusional morphology
    in Marathi.
    \bleu{}'s surface-form $n$-gram matching systematically
    underestimates quality for these languages;
    \chrf{} and \bert{}Score are more appropriate.

  \item \textbf{Translation evaluated on 100 sentences only.}
    The translation test set is limited to 100 sentences due
    to inference time constraints.
    A larger evaluation set would provide more statistically
    reliable estimates.
\end{itemize}

% ── 8. CONCLUSION ─────────────────────────────────────────────
\section{Conclusion}
\label{sec:conclusion}

We have presented a systematic study of quadruplet sentence
alignment for the under-resourced Marathi--Kannada language pair,
with English and Hindi as auxiliary languages.
Our key finding is that fine-tuning on quadruplet structures
provides consistent gains over both the zero-shot baseline and
triplet training for English-centric language pairs:
\acc{} improves from 7.8\% to 40.7\% for EN$\leftrightarrow$KN
and from 41.0\% to 46.3\% for EN$\leftrightarrow$MR.

For the direct Marathi$\leftrightarrow$Kannada pair --- the
most practically important and most under-resourced direction ---
the triplet model achieves 30.2\% \acc{} from a 6.8\% baseline,
a 4.5$\times$ improvement.
Translation quality, measured by \bert{}Score F1 ($\sim$0.786--0.808)
and round-trip cosine similarity (0.898), confirms that the English
pivot preserves semantic content adequately for practical use cases.

\paragraph{Future work.}
Immediate extensions include:
(i) training on the full 179,000+ pair corpus with GPU acceleration;
(ii) replacing CosineSimilarityLoss with
  MultipleNegativesRankingLoss to prevent embedding collapse;
(iii) fine-tuning IndicTrans2 on the legal subdomain corpus;
(iv) human evaluation by bilingual legal professionals; and
(v) extending to Tamil and Telugu to create a broader South
Indian legal MT system.

% ── ACKNOWLEDGEMENTS ──────────────────────────────────────────
\section*{Acknowledgements}
The author thanks Vikram Sir for research supervision and
guidance throughout this project.
This work uses the Samanantar dataset~\cite{ramesh2021samanantar},
PMIndia corpus, OPUS-100~\cite{zhang2020opus100},
MILPaC~\cite{milpac2022},
IndicTrans2~\cite{gala2023indictrans2},
and NLLB-200~\cite{nllb2022}.
Experiments were conducted on a MacBook Pro (Apple M-series CPU)
and Google Colab (T4 GPU).

% ── REFERENCES ────────────────────────────────────────────────
\bibliographystyle{plainnat}
\bibliography{references}

% NOTE: Since you may not have a .bib file ready,
% the references below are provided inline as a fallback.
% In Overleaf, delete the \bibliography line above and
% uncomment the thebibliography block below, OR
% create a references.bib file with the entries provided
% at the end of this file.

\end{document}
